\chapter{Error boundaries and recovery}
\label{ch:errorboundary}

\begin{aside}
The Elm architecture treats expected failures as messages. But
what about unexpected ones --- a stray null deref, an assertion
fired, a third-party library throwing an exception you did not
plan for? \maya{} provides \emph{error boundaries}: a way to
isolate a subtree of the UI so that when it fails, only that
subtree breaks, not the whole app. This chapter explains the
mechanism and how to use it.
\end{aside}

\section{What an error boundary is}

An error boundary is a wrapper component that catches exceptions
thrown by its child's \code{view} or \code{update}. When a
catch happens, the boundary swaps in a fallback view (typically a
\emph{something went wrong} message with a retry button) and
isolates the failure.

The pattern comes from React's \emph{error boundary} concept.
The mechanics in \maya{} are different (no class hierarchy, no
\code{componentDidCatch}) but the goal is the same: contain
failure, do not crash the whole UI.

\subsection*{Without an error boundary}

If \code{view} throws, the runtime catches the exception, prints
a diagnostic, restores the terminal, and exits. Your user sees
their app vanish. The reproducer is buried in stdout.

\subsection*{With an error boundary}

If \code{view} of a sub-tree throws, the boundary catches and
renders a small ``error'' panel in place of that sub-tree. The
rest of the UI continues. The user can keep working; the developer
can inspect the error message.

\section{The ErrorBoundary type}

\file{maya/include/maya/app/error\_boundary.hpp} defines:

\begin{lstlisting}
template <ElementLike Child>
class ErrorBoundary {
public:
    using Fallback = std::function<Element(std::exception_ptr)>;

    ErrorBoundary(Child child, Fallback fallback);

    Element build() &&;   // catches exceptions inside child.build()
};

template <ElementLike Child, class F>
auto error_boundary(Child&& child, F&& fallback);
\end{lstlisting}

The wrapper takes a child and a fallback. \code{build} wraps the
child's own \code{build} in a \code{try/catch} block:

\begin{lstlisting}
template <ElementLike Child>
Element ErrorBoundary<Child>::build() && {
    try {
        return std::move(child_).build();
    } catch (...) {
        return fallback_(std::current_exception());
    }
}
\end{lstlisting}

If the child's build (which recursively builds the entire
sub-tree) succeeds, the wrapper passes it through. If it throws
anywhere, the fallback is rendered instead.

\subsection*{The fallback function}

The fallback receives a \code{std::exception\_ptr} so it can
extract a message if desired:

\begin{lstlisting}
auto safe_panel = error_boundary(
    risky_panel(state),
    [](std::exception_ptr ep) {
        std::string msg;
        try { std::rethrow_exception(ep); }
        catch (const std::exception& e) { msg = e.what(); }
        catch (...)                       { msg = \"unknown error\"; }
        return box(
            text(\"This panel crashed:\") | Bold | Fg<255, 120, 120>,
            text(msg) | Dim
        ) | border<Single> | pad<1>;
    }
);
\end{lstlisting}

A user-friendly fallback hides the exception details by default
and reveals them on demand. A developer-mode fallback dumps the
full \code{what()} immediately.

\section{Where boundaries belong}

Not every sub-tree needs a boundary. Adding boundaries
everywhere defeats the purpose (the runtime becomes lossy). Add
one where:

\begin{itemize}[leftmargin=*]
  \item The sub-tree renders user-supplied content (a markdown
    preview, a code highlighter) where input may be malformed.
  \item The sub-tree calls third-party code you do not control.
  \item The sub-tree is independent of the rest of the UI ---
    losing it does not lose work in progress.
\end{itemize}

A reasonable rule: the panels of a multi-panel app each get a
boundary. The whole-app root does not (because if rendering the
whole app fails, there is nothing useful to show).

\section{Catching during update}

\code{ErrorBoundary} only catches in \code{build}/\code{view}.
Exceptions in \code{update} are different: they corrupt the
model.

\maya{} handles update exceptions at the runtime level:

\begin{lstlisting}
try {
    auto [m2, c2] = P::update(model, msg);
    model = std::move(m2);
    interpret(c2, loop);
} catch (...) {
    log_update_exception(std::current_exception(), model, msg);
    // The model is unchanged from before update was called.
    // No effects were interpreted.
}
\end{lstlisting}

If \code{update} throws, the model is not mutated and no
commands are issued. The next frame renders the previous model.
The user might not notice anything happened; the log records it.

This is intentionally conservative: an exception in \code{update}
indicates a programming bug, not an expected condition. The
right fix is for the developer to add a branch in \code{update}
that returns a message, not for the runtime to silently retry.

\begin{warning}
Do not rely on the runtime's update-exception catching as a
defence in depth. It exists to keep the program alive while the
bug is diagnosed, not as a feature.
\end{warning}

\section{Effect failures}

Failures inside \code{Cmd::task} are wrapped at the boundary:

\begin{lstlisting}
return Cmd<Msg>::task([url]() -> Msg {
    try {
        return FetchOk{do_fetch(url)};
    } catch (const std::exception& e) {
        return FetchFailed{e.what()};
    }
});
\end{lstlisting}

The task callable catches its own exceptions and routes them
through the message type. \code{update} branches on the result.
No exception ever crosses the task boundary.

If you fail to wrap, \maya{}'s task scheduler catches the
exception and logs it, but the corresponding message is never
delivered. Your \code{update} will not see a result. This is a
silent failure --- bad enough that you should always wrap.

\section{Recovery patterns}

\subsection*{Retry on a timer}

After a failure, schedule a retry:

\begin{lstlisting}
[&](FetchFailed e) {
    m.error = e.reason;
    m.retry_count++;
    if (m.retry_count < 3) {
        return std::pair{m,
            Cmd<Msg>::after(1s * (1 << m.retry_count),
                            Retry{})};
    }
    return std::pair{m, Cmd<Msg>{}};
}
\end{lstlisting}

Exponential backoff in three lines.

\subsection*{Reset on user action}

Provide a ``retry'' button in the error UI:

\begin{lstlisting}
auto error_panel = v(
    text(error_message) | Fg<255, 120, 120>,
    text(\"[r] retry  [q] quit\") | Dim
);
\end{lstlisting}

Pressing \code{r} fires a \code{Retry} message; \code{update}
re-issues the original \code{Cmd}.

\subsection*{Degraded mode}

If a feature fails repeatedly, disable it and keep going:

\begin{lstlisting}
if (m.feature_x_failures >= 3) {
    m.feature_x_disabled = true;
}
\end{lstlisting}

The view checks the flag and renders a placeholder. The rest of
the app stays alive.

\section{Logging the failure}

A boundary that catches but never reports is worse than no
boundary at all (you do not even know there was a bug). Pair the
boundary with a log:

\begin{lstlisting}
auto safe_panel = error_boundary(
    risky_panel(state),
    [&](std::exception_ptr ep) {
        log_exception(ep, \"risky_panel\");
        return fallback_ui(ep);
    }
);
\end{lstlisting}

Where \code{log\_exception} writes to a file or sends to a
telemetry service. For development, \code{std::cerr} is fine.
For production, write to a debug ring buffer that the user can
toggle on demand.

\begin{aside}
A debug ring buffer is a small in-memory log that survives
crashes. \maya{}'s built-in one keeps the last 1024 messages and
exposes them via a \code{maya::dump\_log()} call --- useful when
the user reports ``something looked weird''.
\end{aside}

\section{Testing error paths}

The boundary itself is testable:

\begin{lstlisting}
TEST_CASE(\"boundary catches throwing build\") {
    auto throwing = test::throw_on_build();
    auto wrapped = error_boundary(throwing, [](auto) {
        return text(\"fallback\");
    });
    Element e = std::move(wrapped).build();
    REQUIRE(is_text(e, \"fallback\"));
}
\end{lstlisting}

A test helper that throws on \code{build}, wrapped in a
boundary; the result should be the fallback. The runtime is
not involved.

\section{Pitfalls}

\begin{warning}
\textbf{Catching too broadly.} A boundary that catches all
exceptions can mask bugs you want to surface. Catch specific
exception types where possible, and re-throw the ones you do not
understand.
\end{warning}

\begin{warning}
\textbf{Side effects in the fallback.} The fallback is a render
function; it should not have side effects (file I/O, network).
If you need to log, do it before returning the fallback element.
\end{warning}

\begin{warning}
\textbf{Boundary on a leaf.} Wrapping a single text node in a
boundary is wasteful (text nodes do not throw). Boundaries pay
off on subtrees that aggregate work.
\end{warning}

\begin{warning}
\textbf{Boundary that itself throws.} If your fallback throws,
the runtime catches it and you crash anyway. Keep fallbacks
defensively simple.
\end{warning}

\section*{Workshop}

\begin{enumerate}[leftmargin=*]
  \item Wrap your app's largest panel in an
    \code{error\_boundary}. Force an exception inside it and
    verify only that panel is replaced.
  \item Add a debug overlay that shows the last logged
    exception when the user presses \code{F12}.
  \item Read
    \file{maya/include/maya/app/error\_boundary.hpp}. Identify
    where the exception is caught and how the fallback is
    invoked.
\end{enumerate}

\begin{recap}
Error boundaries isolate failure. A boundary catches exceptions
during \code{build}/\code{view} and renders a fallback in place
of the broken sub-tree. \code{update} exceptions are caught at
the runtime level and the model is left unchanged. Task
exceptions must be converted to messages at the task boundary.
Together these mechanisms keep \maya{} apps alive in the face of
bugs.
\end{recap}
